\documentclass[12pt,a4paper]{article}
\usepackage[utf8]{inputenc}
\usepackage[T5]{fontenc}
\usepackage[vietnamese]{babel}
\usepackage[margin=1.5cm,top=1.4cm,bottom=1.4cm]{geometry}
\usepackage{amsmath,amssymb,array,tabularx,enumitem,fancyhdr,tikz,xcolor}
\usetikzlibrary{calc,angles,quotes,arrows.meta,babel,patterns}
\definecolor{ink}{RGB}{30,55,70}

\newcommand{\headerQuote}{%
  \ifcase\value{page}\relax
    \or Math gives us hope that every problem has a solution% Trang 1
    \or Math + You = Champion% Trang 2
    \or Where Frustration Becomes Inspiration% Trang 3
    \or Turning Confusion into Clarity% Trang 4
    \or Don't be an X, be a Y of mathematics% Trang 5
    \or Don't get mad. Learn to add% Trang 6
    \or Math gives us hope that every problem has a solution% Trang 7
    \or Math + You = Champion% Trang 8
    \or Where Frustration Becomes Inspiration% Trang 9
    \or Turning Confusion into Clarity% Trang 10
    \or Don't be an X, be a Y of mathematics% Trang 11
    \or Don't get mad. Learn to add% Trang 12
    \else Math gives us hope that every problem has a solution%
  \fi
}
\pagestyle{fancy}\fancyhf{}
\lhead{\small\textit{\headerQuote}}
\rhead{\small\textit{Người biên soạn: Dương Gia Bảo}}
\rfoot{\small Trang \thepage}
\renewcommand{\headrulewidth}{0.4pt}
\setlength{\headheight}{18pt}
\setlength{\parindent}{0pt}
\setlength{\parskip}{1.5pt}
\setlength{\tabcolsep}{4pt}
\renewcommand{\arraystretch}{1.12}
\newcommand{\parttitle}[1]{\par\vspace{2pt}{\large\bfseries\color{ink}#1}\par\vspace{1pt}}
\newcommand{\four}[4]{\par\noindent\begin{tabularx}{\linewidth}{@{}*{4}{>{\raggedright\arraybackslash}X}@{}}\textbf{A.} #1&\textbf{B.} #2&\textbf{C.} #3&\textbf{D.} #4\end{tabularx}}
\newcommand{\two}[4]{\par\noindent\begin{tabularx}{\linewidth}{@{}*{2}{>{\raggedright\arraybackslash}X}@{}}\textbf{A.} #1&\textbf{B.} #2\\[2pt]\textbf{C.} #3&\textbf{D.} #4\end{tabularx}}
\newcommand{\answerbox}{\hfill\textit{Đáp số:}\ \fbox{\rule{0pt}{12pt}\hspace{2.2cm}}}
\setlist[enumerate,1]{label=\textbf{Câu \arabic*.},leftmargin=1.4cm,labelsep=.15cm,itemsep=1.5pt,topsep=1.0pt,parsep=0pt}
\setlist[enumerate,2]{label=\alph*),leftmargin=.6cm,itemsep=1.0pt,topsep=1.0pt,parsep=0pt}

\begin{document}
\noindent
\begin{minipage}[c]{0.38\textwidth}
    \centering
    \textbf{LÀM TOÁN ĐỂ HẠNH PHÚC}\\[-2pt]
    \rule{3.2cm}{0.6pt}\\[4pt]
    \textbf{ĐỀ CHÍNH THỨC}
\end{minipage}\hfill
\begin{minipage}[c]{0.35\textwidth}
    \centering
    \textbf{ĐỀ ÔN THI GK1}\\
    \textbf{NĂM HỌC 2026 -- 2027}\\[2pt]
    \textbf{ĐỀ SỐ 20}\\
    \textbf{Môn: TOÁN -- Khối lớp 10}\\
    \textit{\small Thời gian làm bài: 90 phút\\(không kể thời gian phát đề)}
\end{minipage}

\vspace{0.2cm}
\noindent
\begin{minipage}[b]{0.78\textwidth}
    \textbf{Họ và tên:}\dotfill\textbf{SBD:}\makebox[3.5cm]{\dotfill}
\end{minipage}\hfill
\begin{minipage}[b]{0.20\textwidth}
    \hfill
    \setlength{\fboxrule}{0.8pt}%
    \fbox{\begin{minipage}{2.3cm}
        \centering
        \textbf{\small Mã đề thi}\\
        \textbf{0020}
    \end{minipage}}
\end{minipage}
\vspace{0.2cm}

\parttitle{PHẦN I. Trắc nghiệm nhiều lựa chọn}
\textit{Thí sinh trả lời từ câu 1 đến câu 12. Mỗi câu hỏi chỉ chọn một phương án đúng duy nhất.}
\begin{enumerate}
\item Cho mệnh đề $P: ``\forall x \in \mathbb{R},\; x^2 + 1 > 0"$. Mệnh đề phủ định của mệnh đề $P$ là
\four{$\overline{P}: ``\exists x \in \mathbb{R},\; x^2 + 1 \le 0"$.}{$\overline{P}: ``\forall x \in \mathbb{R},\; x^2 + 1 \le 0"$.}{$\overline{P}: ``\exists x \in \mathbb{R},\; x^2 + 1 < 0"$.}{$\overline{P}: ``\forall x \in \mathbb{R},\; x^2 + 1 < 0"$.}

\item Cho mệnh đề chứa biến $P(x): ``2x^2 - 5x + 2 = 0"$. Trong các giá trị sau, giá trị nào của $x$ làm cho $P(x)$ là mệnh đề \textbf{đúng}?
\four{$x = 1$.}{$x = 2$.}{$x = 0$.}{$x = 3$.}

\item Tập hợp các số thực $x$ thỏa mãn $-3 \le x < 4$ được biểu diễn dưới dạng khoảng, đoạn, nửa khoảng là
\four{$(-3; 4]$.}{$[-3; 4]$.}{$[-3; 4)$.}{$(-3; 4)$.}

\item Cho định lý: ``Nếu hai tam giác bằng nhau thì diện tích của chúng bằng nhau''. Phát biểu nào sau đây sử dụng thuật ngữ ``điều kiện đủ'' là \textbf{đúng}?
\two{Diện tích hai tam giác bằng nhau là điều kiện đủ để hai tam giác đó bằng nhau.}{Hai tam giác có diện tích bằng nhau là điều kiện cần và đủ để hai tam giác đó bằng nhau.}{Hai tam giác bằng nhau là điều kiện cần để diện tích của chúng bằng nhau.}{Hai tam giác bằng nhau là điều kiện đủ để diện tích của chúng bằng nhau.}

\item Đường thẳng bờ của miền nghiệm của bất phương trình $2x - 3y + 6 \le 0$ có phương trình là
\four{$2x - 3y + 6 = 0$.}{$2x + 3y + 6 = 0$.}{$3x - 2y + 6 = 0$.}{$2x - 3y - 6 = 0$.}

\item Cặp số $(x; y)$ nào sau đây là một nghiệm của hệ bất phương trình $\begin{cases} 2x - y + 1 > 0 \\ x + 3y - 6 \le 0 \end{cases}$?
\four{$(2; 2)$.}{$(1; 1)$.}{$(0; 3)$.}{$(3; 2)$.}

\item Hệ bất phương trình nào sau đây là một hệ bất phương trình bậc nhất hai ẩn?
\four{$\begin{cases} x^2 + y \le 2 \\ x - y > 1 \end{cases}$}{$\begin{cases} x + y - z \ge 0 \\ 2x - y < 3 \end{cases}$}{$\begin{cases} 3x - y + 2 \ge 0 \\ x + 2y - 4 < 0 \end{cases}$}{$\begin{cases} xy + 1 > 0 \\ x - 2y \le 3 \end{cases}$}

\item Bất phương trình nào sau đây là bất phương trình bậc nhất hai ẩn?
\four{$2x^2 + 3y \le 0$.}{$x + y^2 > 1$.}{$2x - 3y + z \ge 4$.}{$3x - 5y + 1 \ge 0$.}

\item Cho góc $\alpha$ thỏa mãn $\tan \alpha = 2$. Giá trị của biểu thức $P = \dfrac{3\sin \alpha - 2\cos \alpha}{\sin \alpha + \cos \alpha}$ bằng
\four{$\dfrac{4}{3}$.}{$\dfrac{1}{3}$.}{$2$.}{$\dfrac{5}{3}$.}

\item Cho tam giác $ABC$ có các cạnh $BC = a, CA = b, AB = c$ và bán kính đường tròn ngoại tiếp là $R$. Khẳng định nào sau đây là \textbf{đúng}?
\four{$\dfrac{a}{\sin A} = R$.}{$\dfrac{a}{\sin A} = 2R$.}{$\dfrac{a}{\sin A} = 4R$.}{$\dfrac{a}{\sin A} = \dfrac{R}{2}$.}

\item Cho tam giác $ABC$ có cạnh $b = 5\text{ cm}$, $c = 8\text{ cm}$ và góc $\widehat{A} = 60^\circ$. Diện tích của tam giác $ABC$ bằng
\four{$20\text{ cm}^2$.}{$20\sqrt{2}\text{ cm}^2$.}{$10\sqrt{3}\text{ cm}^2$.}{$10\text{ cm}^2$.}

\item Giá trị của $\sin 150^\circ$ bằng
\four{$-\dfrac{1}{2}$.}{$-\dfrac{\sqrt{3}}{2}$.}{$\dfrac{\sqrt{3}}{2}$.}{$\dfrac{1}{2}$.}
\end{enumerate}

\parttitle{PHẦN II. Trắc nghiệm đúng/sai}
\textit{Thí sinh trả lời từ câu 13 đến câu 14. Trong mỗi ý a), b), c), d) ở mỗi câu, thí sinh chọn đúng hoặc sai.}
\begin{enumerate}[resume]
\item Xét tính đúng sai của các khẳng định sau:
\begin{enumerate}
    \item Số $2026$ là một số chia hết cho $2$.
    \item Tứ giác có hai đường chéo cắt nhau tại trung điểm của mỗi đường là hình bình hành.
    \item Với mọi số thực $x$, ta luôn có $x^2 - 4x + 5 > 0$.
    \item Tồn tại số nguyên $n$ để $n^2 + 3n + 1$ là một số chẵn.
\end{enumerate}

\item Cho góc $\alpha$ thỏa mãn $90^\circ < \alpha < 180^\circ$ và $\sin \alpha = \dfrac{3}{5}$. Xét tính đúng sai của các khẳng định sau:
\begin{enumerate}
    \item Vì $90^\circ < \alpha < 180^\circ$ nên $\cos \alpha < 0$.
    \item Giá trị của $\cos \alpha$ bằng $-\dfrac{4}{5}$.
    \item Giá trị của $\tan \alpha$ bằng $\dfrac{3}{4}$.
    \item Giá trị của biểu thức $P = \dfrac{\sin \alpha + 2\cos \alpha}{\tan \alpha - \cot \alpha}$ bằng $-\dfrac{12}{7}$.
\end{enumerate}
\end{enumerate}

\parttitle{PHẦN III. Trắc nghiệm trả lời ngắn}
\textit{Thí sinh trả lời từ câu 15 đến câu 18. Viết kết quả vào ô quy định.}
\begin{enumerate}[resume]
\item Cho góc $\alpha$ thỏa mãn $90^\circ < \alpha < 180^\circ$ và $\cos \alpha = -\dfrac{5}{13}$.
Tính giá trị của biểu thức $P = 26\sin \alpha - 5\tan \alpha$.
\answerbox

\item Một lớp học có $45$ học sinh. Qua khảo sát sở thích môn học, có $25$ học sinh thích môn Toán, $20$ học sinh thích môn Ngữ văn và $18$ học sinh thích môn Tiếng Anh. Trong đó có $12$ học sinh thích cả Toán và Ngữ văn, $8$ học sinh thích cả Ngữ văn và Tiếng Anh, $10$ học sinh thích cả Toán và Tiếng Anh; có $5$ học sinh thích cả ba môn trên. Hỏi lớp học đó có bao nhiêu học sinh không thích môn nào trong ba môn trên?
\answerbox

\item Cho hệ bất phương trình bậc nhất hai ẩn $\begin{cases} x \ge 0 \\ y \ge 0 \\ x + y \le 4 \\ 2x + y \le 6 \end{cases}$. Miền nghiệm của hệ bất phương trình chứa bao nhiêu điểm $M(x; y)$ có tọa độ là các số nguyên dương?
\answerbox

\item Bác Nam dự định trồng ngô và khoai lang trên diện tích đất $8\text{ ha}$. Mỗi ha trồng ngô cần $10$ ngày công và thu lãi $20$ triệu đồng; mỗi ha trồng khoai lang cần $15$ ngày công và thu lãi $25$ triệu đồng. Biết rằng tổng số ngày công gia đình có không quá $90$ ngày công. Tính số tiền lãi lớn nhất mà bác Nam có thể thu được (đơn vị: triệu đồng).
\answerbox
\end{enumerate}

\parttitle{PHẦN IV. Tự luận}
\textit{Thí sinh trình bày lời giải chi tiết từ câu 19 đến câu 22.}
\begin{enumerate}[resume]
\item \textit{(1,0 điểm)} \begin{minipage}[t]{0.60\linewidth}
Một bức tranh phong cảnh có chiều cao $AB = 1{,}6\text{ m}$ được treo thẳng đứng trên một bức tường phẳng. Một người quan sát đứng ngắm bức tranh sao cho tầm mắt ở vị trí điểm $O$, cách mặt đất một khoảng cố định và thẳng hàng với chân tường tại điểm $H$ ($OH \perp HB$). Biết mép dưới $A$ của bức tranh cách tầm mắt người xem một khoảng $AH = 0{,}9\text{ m}$ (như hình vẽ bên).
\begin{enumerate}
    \item Gọi $x$ (mét, $x > 0$) là khoảng cách $OH$ từ người quan sát đến tường. Biểu diễn $\tan \widehat{AOB}$ theo $x$.
    \item Người quan sát cần đứng cách bức tường một khoảng $x$ bằng bao nhiêu mét để góc nhìn bức tranh $\widehat{AOB}$ đạt giá trị lớn nhất?
\end{enumerate}
\end{minipage}\hfill
\begin{minipage}[t]{0.37\linewidth}
\centering
\begin{tikzpicture}[scale=0.9, every node/.style={scale=0.85}]
    % Tường thẳng đứng
    \draw[line width=1.4pt, gray!60!black] (3.0,-0.4) -- (3.0,3.6) node[above, font=\bfseries] {Tường};
    % Nền phương ngang
    \draw[dashed, gray] (-0.2,0) -- (3.0,0);
    
    % Các điểm
    \coordinate (O) at (0,0);
    \coordinate (H) at (3.0,0);
    \coordinate (A) at (3.0,0.9);
    \coordinate (B) at (3.0,2.5);

    % Bức tranh AB
    \draw[line width=3.5pt, orange!90!black] (A) -- (B);
    \node[right, font=\bfseries, orange!90!black] at (3.0,1.7) {Bức tranh};

    % Tia ngắm
    \draw[thick, ink] (O) node[left] {$O\text{ (mắt)}$} -- (H) node[below right] {$H$};
    \draw[thick, blue!80!black] (O) -- (A) node[above right] {$A$};
    \draw[thick, teal!80!black] (O) -- (B) node[above right] {$B$};

    % Ký hiệu góc vuông tại H
    \draw (2.7,0) -- (2.7,0.3) -- (3.0,0.3);

    % Chấm điểm
    \filldraw[black] (O) circle (2pt);
    \filldraw[black] (H) circle (2pt);
    \filldraw[blue!80!black] (A) circle (2pt);
    \filldraw[teal!80!black] (B) circle (2pt);

    % Nhãn kích thước
    \node[below] at (1.5,0) {$x = OH$};
    \draw[<->, >=Stealth, red!80!black] (3.5,0) -- (3.5,0.9) node[midway, right] {$0{,}9\text{ m}$};
    \draw[<->, >=Stealth, red!80!black] (3.5,0.9) -- (3.5,2.5) node[midway, right] {$1{,}6\text{ m}$};
\end{tikzpicture}
\end{minipage}

\item \textit{(0,75 điểm)} Cho tam giác $ABC$ có các cạnh $BC = a, CA = b, AB = c$ và diện tích $S$.
\begin{enumerate}
    \item Chứng minh rằng: $S = \dfrac{1}{4}(a^2 + b^2 - c^2)\tan C$.
    \item Chứng minh rằng nếu các góc của tam giác thỏa mãn hệ thức $\sin A = 2\sin B \cos C$ thì tam giác $ABC$ là một tam giác cân tại đỉnh $A$.
\end{enumerate}

\item \textit{(0,75 điểm)} Cho hai tập hợp số $A = [-3; 2)$ và $B = (0; 5]$.
\begin{enumerate}
    \item Xác định các tập hợp $A \cap B$, $A \cup B$, $A \setminus B$ và $B \setminus A$.
    \item Biểu diễn tập hợp $A \cap B$ và $A \cup B$ trên trục số.
\end{enumerate}

\item \textit{(0,5 điểm)} Trong mặt phẳng tọa độ $Oxy$:
\begin{enumerate}
    \item Vẽ đường thẳng $d: 2x - y - 2 = 0$.
    \item Biểu diễn hình học miền nghiệm của bất phương trình $2x - y - 2 \le 0$ trên mặt phẳng tọa độ $Oxy$ (gạch bỏ phần nửa mặt phẳng không thuộc miền nghiệm và nêu rõ kết luận).
\end{enumerate}
\end{enumerate}

\newpage
% =========================================================================
% BẢNG ĐÁP ÁN VÀ LỜI GIẢI CHI TIẾT
% =========================================================================
\parttitle{BẢNG ĐÁP ÁN TỔNG HỢP}

\vspace{0.2cm}
\textbf{1. Phần I: Trắc nghiệm nhiều lựa chọn}
\begin{center}
{\setlength{\tabcolsep}{2.5pt}\small
\begin{tabular}{|c|c|c|c|c|c|c|c|c|c|c|c|}
\hline
\textbf{Câu 1} & \textbf{Câu 2} & \textbf{Câu 3} & \textbf{Câu 4} & \textbf{Câu 5} & \textbf{Câu 6} & \textbf{Câu 7} & \textbf{Câu 8} & \textbf{Câu 9} & \textbf{Câu 10} & \textbf{Câu 11} & \textbf{Câu 12} \\
\hline
\textbf{A} & \textbf{B} & \textbf{C} & \textbf{D} & \textbf{A} & \textbf{B} & \textbf{C} & \textbf{D} & \textbf{A} & \textbf{B} & \textbf{C} & \textbf{D} \\
\hline
\end{tabular}}
\end{center}

\vspace{0.2cm}
\textbf{2. Phần II: Trắc nghiệm đúng/sai}
\begin{center}
\begin{tabular}{|c|c|c|c|c|}
\hline
\textbf{Câu hỏi} & \textbf{Ý a)} & \textbf{Ý b)} & \textbf{Ý c)} & \textbf{Ý d)} \\
\hline
\textbf{Câu 13} & \textbf{Đúng} & \textbf{Đúng} & \textbf{Đúng} & \textbf{Sai} \\
\hline
\textbf{Câu 14} & \textbf{Đúng} & \textbf{Đúng} & \textbf{Sai} & \textbf{Đúng} \\
\hline
\end{tabular}
\end{center}

\vspace{0.2cm}
\textbf{3. Phần III: Trắc nghiệm trả lời ngắn}
\begin{center}
\begin{tabular}{|c|c|c|c|}
\hline
\textbf{Câu 15} & \textbf{Câu 16} & \textbf{Câu 17} & \textbf{Câu 18} \\
\hline
\textbf{36} & \textbf{7} & \textbf{5} & \textbf{170} \\
\hline
\end{tabular}
\end{center}

\vspace{0.2cm}
\textbf{4. Phần IV: Tự luận}
\begin{center}
\begin{tabularx}{\linewidth}{|c|>{\raggedright\arraybackslash}X|}
\hline
\textbf{Câu hỏi} & \textbf{Kết quả tóm tắt} \\
\hline
\textbf{Câu 19} & a) $\tan \widehat{AOB} = \dfrac{1{,}6x}{x^2 + 2{,}25}$;\quad b) Góc nhìn lớn nhất khi $x = 1{,}5\text{ m}$ (với $\tan \widehat{AOB} = \dfrac{8}{15}$). \\
\hline
\textbf{Câu 20} & a) Dùng định lý côsin $a^2 + b^2 - c^2 = 2ab\cos C \implies S = \dfrac{1}{4}(a^2+b^2-c^2)\tan C$;\newline b) $\sin A = 2\sin B \cos C \iff a = \dfrac{a^2+b^2-c^2}{a} \iff b = c \implies \Delta ABC$ cân tại $A$. \\
\hline
\textbf{Câu 21} & a) $A \cap B = (0; 2)$, $A \cup B = [-3; 5]$, $A \setminus B = [-3; 0]$, $B \setminus A = [2; 5]$;\newline b) Biểu diễn chính xác các khoảng, đoạn trên hai trục số. \\
\hline
\textbf{Câu 22} & a) $d$ qua $(0; -2)$ và $(1; 0)$;\quad b) Miền nghiệm là nửa mặt phẳng bờ $d$ chứa gốc tọa độ $O(0; 0)$ (kể cả bờ $d$). \\
\hline
\end{tabularx}
\end{center}

\vspace{0.4cm}
% =========================================================================
% HÌNH VẼ BIỂU DIỄN MIỀN NGHIỆM (HỆ BẤT PHƯƠNG TRÌNH)
% =========================================================================
\parttitle{HÌNH VẼ BIỂU DIỄN MIỀN NGHIỆM (HỆ BẤT PHƯƠNG TRÌNH)}

\begin{minipage}[t]{0.48\textwidth}
\centering
\textbf{Hình 1: Miền nghiệm Câu 18 (Quy hoạch tuyến tính)}\\[6pt]
\begin{tikzpicture}[scale=0.55, >=Stealth]
    % Trục toạ độ
    \draw[line width=1.2pt, ->] (-1.0,0) -- (10.0,0) node[right, font=\normalsize\bfseries] {$x$};
    \draw[line width=1.2pt, ->] (0,-1.0) -- (0,8.5) node[above, font=\normalsize\bfseries] {$y$};
    \node[below left, font=\normalsize\bfseries] at (0,0) {$O$};

    % Đa giác nghiệm O(0;0), A(8;0), B(6;2), C(0;6)
    \fill[fill=ink!18] (0,0) -- (8,0) -- (6,2) -- (0,6) -- cycle;
    \draw[line width=1.8pt, blue!70!black] (0,0) -- (8,0) -- (6,2) -- (0,6) -- cycle;

    % d1: x + y = 8 (cắt (8;0) và (0;8))
    \draw[line width=1.1pt, dashed, red!80!black] (4.5, 3.5) -- (8.8, -0.8) node[below right, font=\small] {$d_1: x+y=8$};
    % d2: 2x + 3y = 18 (cắt (9;0) và (0;6))
    \draw[line width=1.1pt, dashed, teal!80!black] (-0.8, 6.53) -- (9.8, -0.53) node[above right, font=\small] {$d_2: 2x+3y=18$};

    % Đỉnh
    \filldraw[black] (0,0) circle (3pt);
    \filldraw[black] (8,0) circle (3.5pt) node[below, font=\normalsize\bfseries] {$A(8;0)$};
    \filldraw[black] (6,2) circle (3.5pt) node[above right, font=\normalsize\bfseries] {$B(6;2)$};
    \filldraw[black] (0,6) circle (3.5pt) node[left, font=\normalsize\bfseries] {$C(0;6)$};

    % Dóng tọa độ B
    \draw[dotted, thick] (6,0) node[below, font=\small] {$6$} -- (6,2) -- (0,2) node[left, font=\small] {$2$};
\end{tikzpicture}
\end{minipage}\hfill
\begin{minipage}[t]{0.48\textwidth}
\centering
\textbf{Hình 2: Miền nghiệm Câu 22 ($2x - y - 2 \le 0$)}\\[6pt]
\begin{tikzpicture}[scale=0.75, >=Stealth]
    % Trục toạ độ
    \draw[line width=1.2pt, ->] (-2.5,0) -- (4.0,0) node[right, font=\normalsize\bfseries] {$x$};
    \draw[line width=1.2pt, ->] (0,-3.5) -- (0,3.5) node[above, font=\normalsize\bfseries] {$y$};
    \node[above left, font=\normalsize\bfseries] at (0,0) {$O$};

    % Miền nghiệm nửa mặt phẳng chứa O(0;0)
    % Cắt miền [-2.2, 3.5] x [-3.2, 3.2]
    \fill[fill=ink!18] (-2.2, 3.2) -- (2.6, 3.2) -- (-0.6, -3.2) -- (-2.2, -3.2) -- cycle;

    % Đường thẳng bờ d: 2x - y - 2 = 0 <=> y = 2x - 2
    \draw[line width=1.8pt, teal!80!black] (-0.6, -3.2) -- (2.6, 3.2) node[above right, font=\normalsize\bfseries] {$d: 2x - y - 2 = 0$};

    % Điểm A(0; -2) và B(1; 0)
    \filldraw[black] (0,-2) circle (3.5pt) node[right, font=\normalsize\bfseries] {$A(0;-2)$};
    \filldraw[black] (1,0) circle (3.5pt) node[below right, font=\normalsize\bfseries] {$B(1;0)$};
    \filldraw[black] (0,0) circle (3pt);
\end{tikzpicture}
\end{minipage}

\newpage
% =========================================================================
% HƯỚNG DẪN GIẢI CHI TIẾT
% =========================================================================
\parttitle{HƯỚNG DẪN GIẢI CHI TIẾT}

\textbf{PHẦN I. Trắc nghiệm nhiều lựa chọn}
\begin{itemize}[leftmargin=0.6cm, itemsep=3pt]
    \item \textbf{Câu 1:} Phủ định của mệnh đề chứa ký hiệu với mọi ``$\forall x \in X,\; P(x)$'' là mệnh đề ``$\exists x \in X,\; \overline{P(x)}$''. Phủ định của $x^2 + 1 > 0$ là $x^2 + 1 \le 0$. Do đó $\overline{P}: ``\exists x \in \mathbb{R},\; x^2 + 1 \le 0"$. Chọn \textbf{A}.
    \item \textbf{Câu 2:} Thay $x = 2$ vào phương trình ta được $2(2)^2 - 5(2) + 2 = 8 - 10 + 2 = 0$ (mệnh đề đúng). Chọn \textbf{B}.
    \item \textbf{Câu 3:} Tập hợp các số thực $x$ thỏa mãn $-3 \le x < 4$ lấy đầu mút $-3$ (ngoặc vuông) và không lấy đầu mút $4$ (ngoặc tròn), được ký hiệu là nửa khoảng $[-3; 4)$. Chọn \textbf{C}.
    \item \textbf{Câu 4:} Mệnh đề có dạng $P \Rightarrow Q$: ``Nếu hai tam giác bằng nhau ($P$) thì diện tích của chúng bằng nhau ($Q$)''. Khi đó $P$ là điều kiện đủ để có $Q$, tức ``Hai tam giác bằng nhau là điều kiện đủ để diện tích của chúng bằng nhau''. Chọn \textbf{D}.
    \item \textbf{Câu 5:} Đường thẳng bờ của miền nghiệm của bất phương trình $2x - 3y + 6 \le 0$ là đường thẳng $2x - 3y + 6 = 0$. Chọn \textbf{A}.
    \item \textbf{Câu 6:} Thay tọa độ $(1; 1)$ vào hệ bất phương trình:
    $\begin{cases} 2(1) - 1 + 1 = 2 > 0 \\ 1 + 3(1) - 6 = -2 \le 0 \end{cases}$ (thỏa mãn cả hai bất phương trình). Chọn \textbf{B}.
    \item \textbf{Câu 7:} Hệ bất phương trình bậc nhất hai ẩn chỉ chứa các ẩn $x, y$ với bậc cao nhất là bậc 1 và không có tích các ẩn. Hệ $\begin{cases} 3x - y + 2 \ge 0 \\ x + 2y - 4 < 0 \end{cases}$ thỏa mãn định nghĩa. Chọn \textbf{C}.
    \item \textbf{Câu 8:} Bất phương trình $3x - 5y + 1 \ge 0$ là bất phương trình bậc nhất hai ẩn $x, y$. Các phương án khác chứa $x^2, y^2$ hoặc ẩn thứ ba $z$. Chọn \textbf{D}.
    \item \textbf{Câu 9:} Do $\tan \alpha = 2 \implies \cos \alpha \ne 0$. Chia cả tử và mẫu của biểu thức $P$ cho $\cos \alpha$:
    \[
    P = \dfrac{3\dfrac{\sin \alpha}{\cos \alpha} - 2}{\dfrac{\sin \alpha}{\cos \alpha} + 1} = \dfrac{3\tan \alpha - 2}{\tan \alpha + 1} = \dfrac{3(2) - 2}{2 + 1} = \dfrac{4}{3}.
    \]
    Chọn \textbf{A}.
    \item \textbf{Câu 10:} Theo định lý sin trong tam giác $ABC$, ta có $\dfrac{a}{\sin A} = \dfrac{b}{\sin B} = \dfrac{c}{\sin C} = 2R$. Chọn \textbf{B}.
    \item \textbf{Câu 11:} Diện tích tam giác $ABC$ tính theo công thức:
    \[
    S = \dfrac{1}{2}bc\sin A = \dfrac{1}{2} \cdot 5 \cdot 8 \cdot \sin 60^\circ = 20 \cdot \dfrac{\sqrt{3}}{2} = 10\sqrt{3}\text{ cm}^2.
    \]
    Chọn \textbf{C}.
    \item \textbf{Câu 12:} Sử dụng mối liên hệ giữa hai góc bù nhau: $\sin 150^\circ = \sin(180^\circ - 30^\circ) = \sin 30^\circ = \dfrac{1}{2}$. Chọn \textbf{D}.
\end{itemize}

\vspace{0.2cm}
\textbf{PHẦN II. Trắc nghiệm đúng/sai}
\begin{itemize}[leftmargin=0.6cm, itemsep=3pt]
    \item \textbf{Câu 13:}
    \begin{enumerate}
        \item \textbf{Đúng:} Số $2026$ có chữ số tận cùng là $6$ nên chia hết cho $2$.
        \item \textbf{Đúng:} Đây là dấu hiệu nhận biết hình bình hành đã học trong hình học THCS.
        \item \textbf{Đúng:} Ta có $x^2 - 4x + 5 = (x - 2)^2 + 1$. Vì $(x - 2)^2 \ge 0$ với mọi $x \in \mathbb{R}$ nên $(x - 2)^2 + 1 \ge 1 > 0,\; \forall x \in \mathbb{R}$.
        \item \textbf{Sai:} Với mọi số nguyên $n$, hai số $n$ và $n + 3$ luôn khác tính chẵn lẻ, do đó tích $n(n + 3) = n^2 + 3n$ luôn là một số chẵn. Suy ra $n^2 + 3n + 1$ luôn là một số lẻ với mọi $n \in \mathbb{Z}$. Vậy không tồn tại số nguyên $n$ nào để biểu thức nhận giá trị chẵn.
    \end{enumerate}

    \item \textbf{Câu 14:}
    \begin{enumerate}
        \item \textbf{Đúng:} Với góc tù $90^\circ < \alpha < 180^\circ$, điểm biểu diễn trên nửa đường tròn đơn vị có hoành độ âm nên $\cos \alpha < 0$.
        \item \textbf{Đúng:} $\cos^2 \alpha = 1 - \sin^2 \alpha = 1 - \left(\dfrac{3}{5}\right)^2 = \dfrac{16}{25}$. Vì $\cos \alpha < 0$ nên $\cos \alpha = -\sqrt{\dfrac{16}{25}} = -\dfrac{4}{5}$.
        \item \textbf{Sai:} $\tan \alpha = \dfrac{\sin \alpha}{\cos \alpha} = \dfrac{3/5}{-4/5} = -\dfrac{3}{4} \ne \dfrac{3}{4}$.
        \item \textbf{Đúng:} Ta có $\cot \alpha = \dfrac{1}{\tan \alpha} = -\dfrac{4}{3}$. Khi đó:
        \[
        \text{Tử} = \sin \alpha + 2\cos \alpha = \dfrac{3}{5} + 2\left(-\dfrac{4}{5}\right) = -1.
        \]
        \[
        \text{Mẫu} = \tan \alpha - \cot \alpha = -\dfrac{3}{4} - \left(-\dfrac{4}{3}\right) = -\dfrac{3}{4} + \dfrac{4}{3} = \dfrac{7}{12}.
        \]
        Suy ra $P = \dfrac{-1}{7/12} = -\dfrac{12}{7}$.
    \end{enumerate}
\end{itemize}

\vspace{0.2cm}
\textbf{PHẦN III. Trắc nghiệm trả lời ngắn}
\begin{itemize}[leftmargin=0.6cm, itemsep=3pt]
    \item \textbf{Câu 15:} Vì $90^\circ < \alpha < 180^\circ$ nên $\sin \alpha > 0$.
    Ta có $\sin \alpha = \sqrt{1 - \cos^2 \alpha} = \sqrt{1 - \left(-\dfrac{5}{13}\right)^2} = \sqrt{\dfrac{144}{169}} = \dfrac{12}{13}$.
    Suy ra $\tan \alpha = \dfrac{\sin \alpha}{\cos \alpha} = \dfrac{12/13}{-5/13} = -\dfrac{12}{5}$.
    Giá trị của $P$:
    \[
    P = 26\left(\dfrac{12}{13}\right) - 5\left(-\dfrac{12}{5}\right) = 24 - (-12) = 24 + 12 = 36.
    \]
    \textbf{Đáp số:} $36$.

    \item \textbf{Câu 16:} Gọi $T, V, A$ lần lượt là tập hợp các học sinh thích môn Toán, Ngữ văn và Tiếng Anh.
    Áp dụng công thức số phần tử của hợp ba tập hợp hữu hạn:
    \[
    |T \cup V \cup A| = |T| + |V| + |A| - (|T \cap V| + |V \cap A| + |T \cap A|) + |T \cap V \cap A|
    \]
    \[
    |T \cup V \cup A| = 25 + 20 + 18 - (12 + 8 + 10) + 5 = 63 - 30 + 5 = 38\text{ (học sinh)}.
    \]
    Số học sinh không thích môn nào trong ba môn là: $45 - 38 = 7$ (học sinh).
    \textbf{Đáp số:} $7$.

    \item \textbf{Câu 17:} Điểm $M(x; y)$ có tọa độ là các số nguyên dương nên $x, y \in \mathbb{Z}^+$ ($x \ge 1, y \ge 1$).
    Từ hệ bất phương trình:
    \begin{itemize}
        \item Với $x = 1$: $\begin{cases} y \le 4 - 1 = 3 \\ y \le 6 - 2 = 4 \end{cases} \iff 1 \le y \le 3 \implies y \in \{1; 2; 3\}$ (có $3$ điểm).
        \item Với $x = 2$: $\begin{cases} y \le 4 - 2 = 2 \\ y \le 6 - 4 = 2 \end{cases} \iff 1 \le y \le 2 \implies y \in \{1; 2\}$ (có $2$ điểm).
        \item Với $x = 3$: $\begin{cases} y \le 4 - 3 = 1 \\ y \le 6 - 6 = 0 \end{cases} \implies y \le 0$ (không có $y$ nguyên dương thỏa mãn).
        \item Với $x \ge 4$: $2x + y \ge 2(4) + 1 = 9 > 6$ (không có điểm thỏa mãn).
    \end{itemize}
    Vậy có tất cả $3 + 2 = 5$ điểm nguyên dương thuộc miền nghiệm.
    \textbf{Đáp số:} $5$.

    \item \textbf{Câu 18:} Gọi $x, y$ (ha) lần lượt là diện tích đất dùng để trồng ngô và khoai lang ($x \ge 0, y \ge 0$).
    Theo bài ra ta có hệ ràng buộc:
    \[
    \begin{cases} x \ge 0 \\ y \ge 0 \\ x + y \le 8 \\ 10x + 15y \le 90 \iff 2x + 3y \le 18 \end{cases}
    \]
    Hàm lợi nhuận thu được: $F(x, y) = 20x + 25y$ (triệu đồng).
    Miền nghiệm của hệ là miền tứ giác $OABC$ với tọa độ các đỉnh:
    - $O(0; 0) \implies F(0, 0) = 0$.
    - $A(8; 0) \implies F(8, 0) = 20(8) = 160$ triệu đồng.
    - $B(6; 2)$ (giao điểm của $x + y = 8$ và $2x + 3y = 18$) $\implies F(6, 2) = 20(6) + 25(2) = 170$ triệu đồng.
    - $C(0; 6)$ (giao điểm của $2x + 3y = 18$ với trục $Oy$) $\implies F(0, 6) = 25(6) = 150$ triệu đồng.
    Vậy số tiền lãi lớn nhất thu được là $170$ triệu đồng.
    \textbf{Đáp số:} $170$.
\end{itemize}

\vspace{0.2cm}
\textbf{PHẦN IV. Tự luận}
\begin{itemize}[leftmargin=0.6cm, itemsep=4pt]
    \item \textbf{Câu 19:}
    \begin{enumerate}
        \item Trong tam giác vuông $OHA$ vuông tại $H$: $\tan \widehat{AOH} = \dfrac{AH}{OH} = \dfrac{0{,}9}{x}$.\\
        Độ dài đoạn thẳng $HB = HA + AB = 0{,}9 + 1{,}6 = 2{,}5\text{ m}$.\\
        Trong tam giác vuông $OHB$ vuông tại $H$: $\tan \widehat{BOH} = \dfrac{HB}{OH} = \dfrac{2{,}5}{x}$.\\
        Vì tia $OA$ nằm giữa hai tia $OH$ và $OB$ nên $\widehat{AOB} = \widehat{BOH} - \widehat{AOH}$.\\
        Áp dụng công thức biến đổi tang của một hiệu:
        \[
        \tan \widehat{AOB} = \dfrac{\tan \widehat{BOH} - \tan \widehat{AOH}}{1 + \tan \widehat{BOH} \cdot \tan \widehat{AOH}} = \dfrac{\dfrac{2{,}5}{x} - \dfrac{0{,}9}{x}}{1 + \dfrac{2{,}5}{x} \cdot \dfrac{0{,}9}{x}} = \dfrac{\dfrac{1{,}6}{x}}{\dfrac{x^2 + 2{,}25}{x^2}} = \dfrac{1{,}6x}{x^2 + 2{,}25}.
        \]
        \item Chia cả tử và mẫu cho $x > 0$, ta có:
        \[
        \tan \widehat{AOB} = \dfrac{1{,}6}{x + \dfrac{2{,}25}{x}}.
        \]
        Áp dụng bất đẳng thức Cauchy cho hai số dương $x$ và $\dfrac{2{,}25}{x}$:
        \[
        x + \dfrac{2{,}25}{x} \ge 2\sqrt{x \cdot \dfrac{2{,}25}{x}} = 2\sqrt{2{,}25} = 2 \cdot 1{,}5 = 3.
        \]
        Do đó $\tan \widehat{AOB} \le \dfrac{1{,}6}{3} = \dfrac{8}{15}$.\\
        Vì hàm số tang đồng biến trên khoảng nhọn $(0^\circ; 90^\circ)$ nên góc $\widehat{AOB}$ đạt giá trị lớn nhất khi và chỉ khi $\tan \widehat{AOB}$ lớn nhất, tức bằng $\dfrac{8}{15}$.\\
        Dấu đẳng thức xảy ra khi:
        \[
        x = \dfrac{2{,}25}{x} \iff x^2 = 2{,}25 \iff x = 1{,}5\text{ (m) (do } x > 0\text{)}.
        \]
        Vậy người quan sát cần đứng cách bức tường $1{,}5\text{ m}$ để góc nhìn bức tranh là lớn nhất.
    \end{enumerate}

    \item \textbf{Câu 20:}
    \begin{enumerate}
        \item Theo định lý côsin trong tam giác $ABC$:
        \[
        c^2 = a^2 + b^2 - 2ab\cos C \implies a^2 + b^2 - c^2 = 2ab\cos C.
        \]
        Thay vào vế phải ta được:
        \[
        \text{VP} = \dfrac{1}{4}(a^2 + b^2 - c^2)\tan C = \dfrac{1}{4}(2ab\cos C) \cdot \dfrac{\sin C}{\cos C} = \dfrac{1}{2}ab\sin C = S = \text{VT (đpcm)}.
        \]
        \item Theo định lý sin: $\dfrac{a}{\sin A} = \dfrac{b}{\sin B} = 2R \implies \sin A = \dfrac{a}{2R},\; \sin B = \dfrac{b}{2R}$.\\
        Theo định lý côsin: $\cos C = \dfrac{a^2 + b^2 - c^2}{2ab}$.\\
        Thay vào hệ thức $\sin A = 2\sin B \cos C$, ta được:
        \[
        \dfrac{a}{2R} = 2 \cdot \dfrac{b}{2R} \cdot \dfrac{a^2 + b^2 - c^2}{2ab} \iff a = \dfrac{a^2 + b^2 - c^2}{a} \iff a^2 = a^2 + b^2 - c^2
        \]
        \[
        \iff b^2 - c^2 = 0 \iff b^2 = c^2 \iff b = c\text{ (do } b, c > 0\text{)}.
        \]
        Vì $b = c$ nên tam giác $ABC$ cân tại đỉnh $A$ (đpcm).
    \end{enumerate}

    \item \textbf{Câu 21:}
    \begin{enumerate}
        \item Với $A = [-3; 2)$ và $B = (0; 5]$:
        \begin{itemize}
            \item $A \cap B = (0; 2)$.
            \item $A \cup B = [-3; 5]$.
            \item $A \setminus B = [-3; 0]$.
            \item $B \setminus A = [2; 5]$.
        \end{itemize}
        \item Biểu diễn trên trục số:
        \begin{itemize}
            \item Biểu diễn tập hợp $A \cap B = (0; 2)$: Vẽ trục số $Ox$, xác định hai mốc $0$ và $2$. Dùng dấu ngoặc tròn $($ tại điểm $0$ và dấu ngoặc tròn $)$ tại điểm $2$; gạch bỏ phần bên ngoài khoảng $(0; 2)$.
            \item Biểu diễn tập hợp $A \cup B = [-3; 5]$: Vẽ trục số $Ox$, xác định hai mốc $-3$ và $5$. Dùng dấu ngoặc vuông $[$ tại điểm $-3$ và dấu ngoặc vuông $]$ tại điểm $5$; gạch bỏ phần bên ngoài đoạn $[-3; 5]$.
        \end{itemize}
    \end{enumerate}

    \item \textbf{Câu 22:}
    \begin{enumerate}
        \item Vẽ đường thẳng $d: 2x - y - 2 = 0$:
        \begin{itemize}
            \item Cho $x = 0 \implies y = -2$, ta được điểm $A(0; -2)$.
            \item Cho $y = 0 \implies x = 1$, ta được điểm $B(1; 0)$.
        \end{itemize}
        Vẽ đường thẳng $d$ đi qua hai điểm $A(0; -2)$ và $B(1; 0)$.
        \item Biểu diễn hình học miền nghiệm:
        \begin{itemize}
            \item Lấy điểm gốc tọa độ $O(0; 0) \notin d$. Thay tọa độ điểm $O$ vào vế trái của bất phương trình:
            \[
            2(0) - 0 - 2 = -2 \le 0\text{ (khẳng định đúng)}.
            \]
            \item Vậy miền nghiệm của bất phương trình $2x - y - 2 \le 0$ là nửa mặt phẳng có bờ là đường thẳng $d$, chứa gốc tọa độ $O(0; 0)$ (kể cả bờ $d$).
            \item Trên hình vẽ (như Hình 2 ở phần trên), ta gạch bỏ phần nửa mặt phẳng bờ $d$ không chứa gốc tọa độ $O(0; 0)$.
        \end{itemize}
    \end{enumerate}
\end{itemize}

\end{document}
